\section{Power-Calculation Transparency Table}\label{app:D}

This appendix presents whatever per-iteration power derivation is
feasible at writing time against the framework's pre-pinned
\(\mathrm{MDES} = 0.40\) standard-deviation units of \(Y\) and
\(\mathrm{Power}_{\min} = 0.80\) at one-sided \(\alpha = 0.05\). The
sample-size adjustment under autocorrelated residuals follows the
classical AR(1) effective-sample-size correction
\(N_{\mathrm{eff}} = N \cdot (1 - \rho_1)/(1 + \rho_1)\), where
\(\rho_1\) is the first-order residual autocorrelation; this is the
standard textbook adjustment for the variance of an AR(1) sample mean
and is not specifically attributable to a single canonical reference
(it predates the HAC literature). Project inference uses
Newey-West HAC standard errors \citep{neweywest1987} at fixed lag
truncation per \(\S\ref{sec:5.4}\); we report the realized HAC SE
where it is recorded in the verdict pack so that the implied
effective-sample-size adjustment is auditable from the realized SE
itself rather than from a separately-recorded \(\rho_1\) diagnostic.
Cells where the derivation cannot be honestly produced from the
artifacts available at writing time are flagged \emph{inherited} and
cross-referenced to the supervisor-ask in Section~\ref{sec:7}
(\(\S 7.4\)). We do not fabricate any number we cannot derive from a
stable artifact.

\begin{table}[h]
\centering
\small
\begin{tabular}{@{}p{2.7cm} p{1.3cm} p{1.5cm} p{1.4cm} p{1.6cm} p{4.2cm}@{}}
\toprule
Iteration & \(N\) & \(\hat{\rho}_1\) (resid.) & \(N_{\mathrm{eff}}\) & Implied power at MDES\(=0.40\) & Notes \\
\midrule
Pair D                & 134 & inherited       & inherited       & inherited & \(\hat{\rho}_1\) on the primary residuals not re-derived from a stable artifact at writing time; HAC \(L = 12\) was used in inference but the residual autocorrelation diagnostic is not recorded as a standalone number in the verdict pack. Realized HAC SE \(0.02465\) implies the AR(1) adjustment was material. Inherited from Phase-A.0 Rev-5.3.1; re-derivation tabled for \(\S 7.4\). \\
\addlinespace
dev-AI Stage-1        & 134 & inherited       & inherited       & inherited (FAIL closed) & Iteration closed FAIL 2026-05-06 with realized HAC(\(L = 12\)) SE \(0.0847\); the SE is approximately \(3.4\times\) the Pair~D HAC SE on identical \(N\) and lag panel, indicating substantially larger residual variance and/or autocorrelation at the Section-J-narrow stratum. The realized FAIL is sign-flipped, not underpowered, so the post-hoc power question is moot for the verdict; the SE is informative for the next-iteration prior on Section-M (which would inherit a similar autocorrelation structure). \\
\addlinespace
Section M (candidate next iteration) & 134 (R2 sample) & not yet estimated as primary & --- & not yet estimable as primary & R2 sensitivity arm in dev-AI Stage-1 returned \(\hat{\beta} = +0.455\) at \(t = +4.73\), one-sided \(p = 1.13 \times 10^{-6}\). \emph{Post-hoc identified candidate, NOT pre-registered as a primary in any iteration to date.} A Section-M iteration would require a fresh spec-pin; the tabulated row is informational and does not graduate the R2 finding. \\
\addlinespace
FX-vol-CPI primary    & 947 & inherited       & inherited       & inherited & Pre-registered HAC bandwidth \(L = 4\) at the weekly cadence; the residual autocorrelation diagnostic is recorded indirectly through the gate-aggregation reconciliation but not as a single-number artifact. Inherited from Phase-A.0 Rev-5.3.1; re-derivation tabled for \(\S 7.4\). \\
\addlinespace
Phase-A.0 remittance  & ---  & not applicable  & not applicable  & not applicable & Iteration EXITed at the kill-criterion step before estimation; no power computation is meaningful. \\
\addlinespace
P1 SN18               & 8 events (target) & not applicable & not applicable & 0.80 vs.\ \(d = 1.4\); 0.40 vs.\ \(d = 0.8\); 0.13 vs.\ \(d = 0.5\) & Power profile pinned in the parked spec at \(\alpha_{\mathrm{primary}} = 0.0167\), \(N_{\mathrm{events}} = 8\), \(\mathrm{df} = 7\), two-sided. The pre-pinned profile is reproduced verbatim and is not converted to MDES \(= 0.40\) standard-deviation units, because the Cohen-\(d\) parameterization is the spec's own scale. \\
\bottomrule
\end{tabular}
\caption{Per-iteration power derivation transparency. Cells flagged
\emph{inherited} are not re-derived from the per-iteration realized
\(N\) and residual autocorrelation; the framework currently inherits
the \(N_{\min} = 75\), \(\mathrm{Power}_{\min} = 0.80\),
\(\mathrm{MDES} = 0.40\)-SD invariants from the Phase-A.0 specification
chain. The Section-M row is a post-hoc identified candidate-next-
iteration entry sourced from the dev-AI Stage-1 R2 sensitivity arm and
is \emph{not} a pre-registered primary in any iteration to date.
Re-derivation per iteration is the dedicated supervisor ask in
\(\S 7.4\).}\label{tab:powerD}
\end{table}

The table is the most direct substantive response to the
\(\S 7.4\) ask: it converts the question ``is \(N_{\min} = 75\)
defensible?'' from rhetorical to data-supported wherever the
data are available, and it surfaces honestly the cells where the data
are not available at writing time. The framework does not currently
record a single-number residual-autocorrelation diagnostic in the
verdict pack of either Pair~D, dev-AI Stage-1, or FX-vol-CPI;
remediating that artifact gap is itself one of the methodological
housekeeping tasks that any forward iteration would inherit,
conditional on the supervisor's guidance.

